\setcounter{table}{0}
\renewcommand{\thetable}{C-\arabic{table}}

\section{附录 C　数据集构成与标注规范}

本附录说明项目所用数据集的组成、组织方式与标注规范，作为正文第八章“2.数据支撑”的补充。

\subsection{C.1 数据集组成}

\begin{table}[htbp]
    \centering
    \caption{项目所用数据集的组成与作用}
    \label{tab:ds-compose}
    \small
    \begin{tabular}{|p{3.2cm}|p{2.2cm}|p{3.8cm}|p{5.2cm}|}
        \hline
        \textbf{数据集} & \textbf{类型} & \textbf{影像来源} & \textbf{在本项目中的作用} \\
        \hline
        自建黑土耕地数据集 & 场景数据集 & 高分二号卫星影像、无人机航拍影像 & 面向黑土地场景的模型训练与效果验证，并作为试点应用与系统演示的数据基础 \\
        \hline
        公开基准数据集 & 验证数据集 & 公开遥感影像数据 & 用于模型结构验证与参数调优，使检测方法与同类研究保持可比口径 \\
        \hline
    \end{tabular}
    \par\vspace{2pt}
    {\footnotesize 注：两类数据集在本项目中分工不同——公开基准数据集用于验证模型结构本身是否有效，自建场景数据集用于验证模型在黑土地场景下是否可用。附录 B.4 所列 mF1 指标即在公开基准数据集上测得。}
\end{table}

\subsection{C.2 数据组织方式}

项目影像与标签按四类目录组织，分别存放 A 时相影像、B 时相影像、变化标签与样本清单。每条样本由一对双时相影像与一幅像素级标签构成：两幅影像为同一区域不同时期的观测结果，标签则逐像素标注该区域在该时间段内是否发生变化，变化区域标记为白色、未变化区域标记为黑色。影像在进入模型前统一裁剪为 $256 \times 256$ 像素，以保证输入尺寸一致。

\subsection{C.3 标注规范}

数据标注遵循“采集—配准—标注—核验”的流程，各环节要点如下：

\begin{itemize}
    \item \textbf{影像筛选}：选取时间基线一致、云覆盖较少、成像质量合格的影像对，避免因云影、时相差异过大引入伪变化；
    \item \textbf{几何配准与辐射归一化}：使两个时相的影像在同一空间基准上逐像元对应，并压缩成像条件差异造成的色调偏差；
    \item \textbf{像素级标注}：使用 ArcGIS 等专业工具逐像素勾绘变化区域边界，形成与影像对一一对应的真值标签；
    \item \textbf{交叉复核}：采用双人独立标注并交叉核对的方式控制标注误差，对边界模糊、地类易混淆的样本进行复判或剔除；
    \item \textbf{负样本纳入}：将正常农事活动引起的地表变化（翻耕、收割等）作为负样本纳入数据集，使模型学习“发生了变化但并非退化”的判别边界，降低误判率。
\end{itemize}

上述规范的目的在于保证真值标签的可靠性。变化检测模型的判别能力上限由标签质量决定：若标注存在系统性偏差，模型会学习到错误的判别依据，且在真实场景中难以察觉。因此数据集构建阶段即通过多环节校核把住质量关，为后续模型训练与系统检测提供可信的数据基础。
